\section{AC自動機（Aho-Corasick）}
對於字串問題，我們有一個重型資料結構可以使用，這幾乎是字串問題的大招（荒野亂鬥的極限威能？），那就是 \textbf{AC自動機}（Aho-Corasick Automaton）。

AC自動機不能幫你自動AC題目，但他可以幫你AC很多字串匹配的問題。

由於字串的題目幾乎都可以用之前講過的演算法解決，因此這邊不詳細介紹，僅提供oi-wiki上的程式碼，如果有興趣可以去查看看

\subsection{AC自動機的程式碼}

\lstset{caption={AC自動機 C++ 實作}}
\begin{lstlisting}[language=C++]
#include <cstdio>
#include <cstring>
#include <queue>
using namespace std;

constexpr int N = 2e5 + 6;
constexpr int LEN = 2e6 + 6;
constexpr int SIZE = 2e5 + 6;

int n;

struct Node {
    int son[26];
    int ans;
    int fail;
    int du;
    int idx;

    void init() {
        memset(son, 0, sizeof(son));
        ans = fail = idx = 0;
    }
} tr[SIZE];

int tot;
int ans[N], pidx;

void init() {
    tot = pidx = 0;
    tr[0].init();
}

void insert(char s[], int &idx) {
    int u = 0;
    for (int i = 1; s[i]; i++) {
        int &son = tr[u].son[s[i] - 'a'];
        if (!son) son = ++tot, tr[son].init();
        u = son;
    }
    if (!tr[u].idx) tr[u].idx = ++pidx;
    idx = tr[u].idx;
}

void build() {
    queue<int> q;
    for (int i = 0; i < 26; i++)
        if (tr[0].son[i]) q.push(tr[0].son[i]);
    while (!q.empty()) {
        int u = q.front();
        q.pop();
        for (int i = 0; i < 26; i++) {
            if (tr[u].son[i]) {
                tr[tr[u].son[i]].fail = tr[tr[u].fail].son[i];
                tr[tr[tr[u].fail].son[i]].du++;
                q.push(tr[u].son[i]);
            } else
                tr[u].son[i] =
                    tr[tr[u].fail].son[i];
        }
    }
}

void query(char t[]) {
    int u = 0;
    for (int i = 1; t[i]; i++) {
        u = tr[u].son[t[i] - 'a'];
        tr[u].ans++;
    }
}

void topu() {
    queue<int> q;
    for (int i = 0; i <= tot; i++)
        if (tr[i].du == 0) q.push(i);
    while (!q.empty()) {
        int u = q.front();
        q.pop();
        ans[tr[u].idx] = tr[u].ans;
        int v = tr[u].fail;
        tr[v].ans += tr[u].ans;
        if (!--tr[v].du) q.push(v);
    }
}

char s[LEN];
int idx[N];

int main() {
    AC::init();
    scanf("%d", &n);
    for (int i = 1; i <= n; i++) {
        scanf("%s", s + 1);
        AC::insert(s, idx[i]);
        AC::ans[i] = 0;
    }
    AC::build();
    scanf("%s", s + 1);
    AC::query(s);
    AC::topu();
    for (int i = 1; i <= n; i++) {
        printf("%d\n", AC::ans[idx[i]]);
    }
    return 0;
}

\end{lstlisting}
